\section{Ветвь Мерсенна и арифметический зоопарк}\label{sec:zoo}

За пределами задач Клэя прочтение через двигатель применимо, вскользь, к целому зоопарку
классических вопросов о простых. Этот раздел --- обзор: для каждого семейства мы фиксируем, что
доказано по-настоящему (вложение в язык $6m\pm1$, критерий близнецов, классическая делимость),
чего стоило бы опровержение (фронт манифестации --- зеркало \cref{sec:riemann}) и где
останавливается честность. Объединяющая дисциплина --- намеренная сдержанность: ни для одного из
этих семейств граница декрета \emph{не берётся}. Каждая трилемма пройдена, каждая граница
машинно-допустима, и каждая специально оставлена открытой --- выбор, который мы проговариваем ниже.
Всюду классические факты --- это факты \emph{Начал} Евклида~\cite{euclid-elements} и
Харди--Райта~\cite{hardy-wright}.

\subsection{Ветвь Мерсенна: честный мост}

Числа Мерсенна $2^p-1$ выглядят как готовые обитатели решётки $6m\pm1$, и есть искушение объявить,
что двигатель близнецов решает их «попутно». Это не так. Сначала фиксируем подлинное вложение.

\begin{theorem}[\code{mersenne\_eq\_sixCenter\_add\_one}]\label{thm:mersenne_eq_sixCenter_add_one}
\green{}\; Для нечётного $p$, полагая $m_p=(2^{p-1}-1)/3$,
\begin{equation}\label{eq:mersenne-embed}
  2^p-1 = 6\,m_p+1.
\end{equation}
\end{theorem}

Значит, всякое нечётное число Мерсенна есть в точности плюс-сторона центра $m_p$; его нижняя
сторона $6m_p-1 = 2^p-3$ даёт немедленный критерий близнецов.

\begin{theorem}[\code{isTwinCenter\_mersenneCenter\_iff}]\label{thm:isTwinCenter_mersenneCenter_iff}
\green{}\; Для нечётного $p\ge 2$ центр $m_p$ является центром-близнецом тогда и только тогда, когда
оба числа $2^p-3$ и $2^p-1$ просты.
\end{theorem}

Такие близнецы Мерсенна существуют ($p=3$ даёт $(5,7)$, $p=5$ даёт $(29,31)$), и единственная
честная стрелка между темами исходит \emph{из} них: подпоследовательность близнецов --- всё ещё
близнецы.

\begin{theorem}[\code{twinLowersInfinite\_of\_mersenneTwins}]\label{thm:twinLowersInfinite_of_mersenneTwins}
\green{}\; Если центры-близнецы Мерсенна неограниченны, то \code{TwinLowers.Infinite}.
\end{theorem}

\emph{Идея доказательства.} Тогда пары $(2^p-3,\,2^p-1)$ --- неограниченные пары близнецов, и
подпоследовательность близнецов проносится по честному мосту к подлинной цели программы. Стрелка
тривиальна и ведёт от более сильной гипотезы к более слабому следствию; обратная импликация
«близнецы $\Rightarrow$ Мерсенн» здесь \emph{не} теорема --- ветвь несёт явный маркер покрытия
\code{NoTwinsToMersenneImplicationClaimed}, а цель \code{MersennePrimesInfinite} остаётся \red{}
открытой, ниоткуда в ветви не выводимой.

\begin{remark}
Более ранний «прямой ряд» из $34$ кирпичей стремился доказать бесконечность Мерсенна безусловно;
внутренний аудит вскрыл его как \warn{}~вакуумный --- его пакеты \code{noEngine} были ненаселены, а
заглавные теоремы следовали из пустых посылок. Он зафиксирован, а не заретуширован: цепи редукции
здесь потребляют лишь подлинные гипотезы.
\end{remark}

\subsection{Опровержение предъявляет двигатель, и граница, которую мы не взяли}

Пропустим ветвь Мерсенна через архитектуру манифестации \cref{sec:riemann}. Отклонение здесь ---
не точка-данное, а $\Pi$-свидетель \emph{конечности}: \code{MersenneTwinAbsenceAbove}$(P)$,
утверждающий, что всякая пара простых $(2^p-3,\,2^p-1)$ сидит нижней стороной не выше $P$. Проводка
связывает его с целью точно (\code{mersenneTwinCentersUnbounded\_iff\_no\_absence}, \green{}), а
свидетель нельзя подделать: центры образуют цепь по основанию $4$, $m\to 4m+1$, первый шаг которой
$5\to 1$ не несёт собственного пилинга ни на одном масштабе
(\code{isEmpty\_properCenterPeel\_five\_one}, \green{}), ибо ведущий коэффициент $4$ не прост и не
$\equiv\pm1\pmod 6$ --- стороны не отображаются в стороны. Значит, безусловного запаса потоков здесь
не построить --- в точности риманова ситуация.

\begin{theorem}[\code{mersenneRefutation\_carries\_engine}]\label{thm:mersenneRefutation_carries_engine}
\green{}\; Свидетель отсутствия, закон манифестации \code{MersenneManifestationLaw} и сведённые книги
на масштабе не ниже $P$ предъявляют вечный двигатель как объект,
\code{ConcreteEuclideanEngineWitness}, ещё до всякого убийства.
\end{theorem}

\emph{Идея доказательства.} Разрешающая проекция даёт устойчивую вселенную без энергии; закон
превращает отсутствие в бесконечное семейство потоков; конечный ключ вынужден столкнуть два его
члена; из столкновения собирается свидетель-двигатель. В композиции со стеной «нет двигателей» и
принятой границей это делает близнецов Мерсенна неограниченными, а через мост ---
\code{twinLowersInfinite\_of\_noEngine\_boundary\_and\_mersenneManifestation} (\green{}). Замок по
свидетелю несущий: незапертая форма при $P=0$ дала бы запас на разрешённом масштабе, противореча
зелёной невозможности \code{no\_deviationFlowSupply\_at\_resolved\_scale}.

Трилемма пройдена, и поле границы \code{mersenneBoundary} машинно-допустимо. Его цена точна.

\begin{theorem}[\code{mersenneManifestationLaw\_iff\_unbounded\_of\_boundary}]\label{thm:mersenneManifestationLaw_iff_unbounded_of_boundary}
\green{}\; При принятой границе \code{TheStrictLastStep00Obligation} закон манифестации эквивалентен
неограниченности близнецов Мерсенна:
\begin{equation}\label{eq:mersenne-price}
\begin{aligned}
  &\code{TheStrictLastStep00Obligation}\ \Rightarrow\
  &\quad \bigl(\code{MersenneManifestationLaw}\iff\code{MersenneTwinCentersUnbounded}\bigr).
\end{aligned}
\end{equation}
\end{theorem}

Декрет нёс бы ровно силу близнецов Мерсенна --- и здесь, впервые в программе, знак эвристики
\emph{перевёрнут}: ожидаемое число пар $(2^p-3,\,2^p-1)$ --- сходящаяся сумма, известны лишь
$p=3,5$, а при всяком $p\equiv3\pmod4$ число $2^p-3$ делится на $5$. Взять границу значило бы
поставить против эвристики; мы \emph{намеренно её не взяли}. Машинный критерий говорит «можно»;
честность --- «можно не значит должно».

\subsection{Стороны решётки и узоры Полиньяка: одиночки зелены, пары гипотетичны}

Задачи о парах --- близнецы (зазор $2$), кузены (зазор $4$), sexy (зазор $6$) и Полиньяк вообще ---
раскладываются на решётке геометрически. Сначала хребет честности: каждая сторона по отдельности
несёт бесконечно много простых.

\begin{theorem}[\code{minusSide\_primes\_unbounded}]\label{thm:minusSide_primes_unbounded}
\green{}\; Выше любого $n$ есть простое $p>n$ с $p\bmod 6 = 5$; и, отдельно, простое $p>n$ с
$p\bmod 6 = 1$ (\code{plusSide\_primes\_unbounded}).
\end{theorem}

\emph{Идея доказательства.} Это теорема Дирихле о простых в арифметических прогрессиях, взятая
целиком из \code{mathlib}~\cite{mathlib}: классы $1$ и $5$ по модулю $6$ взаимно просты с $6$,
поэтому каждый несёт бесконечно много простых. Об одном и том же центре два столбца молчат ---
маркер \code{NoPairingClaimed} (\green{}, \code{abbrev := True}) честно это подписывает, а
совпадение сторон на общем центре есть гипотеза о близнецах, здесь не утверждаемая и не выводимая.

Кузены и sexy читаются с той же решётки. Пара кузенов $(p,p+4)$ пересекает два соседних центра,
$(6m+1,\,6(m+1)-1)$; пара sexy $(p,p+6)$ остаётся на одной стороне двух соседних центров. Геометрия
навязывает вычет.

\begin{theorem}[\code{cousin\_mod\_six}]\label{thm:cousin_mod_six}
\green{}\; Всякий младший член пары кузенов с $p>3$ лежит на плюс-стороне: $p\bmod 6 = 1$.
\end{theorem}

\emph{Идея доказательства.} У простого $p>3$ имеем $p\bmod 6\in\{1,5\}$; случай $5$ вынуждает
$p+4\equiv 0\pmod 3$, что невозможно для простого $p+4>3$. Кузены рождаются только на плюс-стороне.

\begin{theorem}[\code{constellation\_of\_adjacent\_twinCenters}]\label{thm:constellation_of_adjacent_twinCenters}
\green{} (даже без \code{Classical.choice})\; Два соседних центра-близнеца $m$ и $m+1$ одновременно
дают кузенов на $m$ и обе пары sexy на $m$.
\end{theorem}

\emph{Идея доказательства.} Тогда все четыре стороны $6m\mp1$, $6(m+1)\mp1$ просты, а кузены
$(6m+1,\,6m+5)$, минус-sexy $(6m-1,\,6m+5)$ и плюс-sexy $(6m+1,\,6m+7)$ суть чистые перестановки ---
при $m=1$ это созвездие $5,7,11,13$. Доказательство не \emph{строит} соседних центров-близнецов; то,
что они существуют сколь угодно высоко, --- гипотеза сильнее близнецов. Оба семейства проходят один
и тот же фронт манифестации, оба поля границ \code{cousinBoundary}, \code{sexyBoundary}
машинно-допустимы, и оба не взяты --- здесь вовсе нет цепи, которую можно пилить, а серийный декрет
стёр бы уникальность, в которой смысл карантина.

\subsection{Софи Жермен: удвоение и классическая делимость}

Пара Софи Жермен --- это $(p,\,2p+1)$, оба простые. В координатах решётки она односторонняя, а
отображение центров SG --- чистое удвоение.

\begin{theorem}[\code{sg\_center\_doubling}]\label{thm:sg_center_doubling}
\green{}\; $2(6m-1)+1 = 6(2m)-1$: у простого SG $p=6m-1$ спутник $2p+1$ лежит на минус-стороне
центра $2m$.
\end{theorem}

Отображение $m\to 2m$ --- третий брат в ряду цепей, рядом с мерсенновой $4m+1$ и пентадической
$5m+1$. Но ветвь несёт единственную содержательную классическую теорему всего зоопарка,
машинно-доказанную до последнего шага.

\begin{theorem}[\code{sophieGermain\_divides\_mersenne}]\label{thm:sophieGermain_divides_mersenne}
\green{}\; Если $p$ просто, $p\equiv 3\pmod 4$, $p>3$ и $q=2p+1$ просто, то $q\mid M_p$ и число
Мерсенна $M_p=2^p-1$ составное.
\end{theorem}

\begin{proof}
Это классический путь Эйлера--Лагранжа~\cite{hardy-wright}. Так как $p\equiv3\pmod4$, имеем
$q=2p+1\equiv7\pmod8$ --- ровно условие $q\equiv\pm1\pmod8$, при котором $2$ есть квадратичный
вычет по модулю $q$ (\code{ZMod.exists\_sq\_eq\_two\_iff}). По критерию Эйлера
$2^{(q-1)/2}\equiv 1\pmod q$, а $(q-1)/2=p$, откуда $2^p\equiv 1\pmod q$, то есть $q\mid 2^p-1=M_p$.
Для $p\ge4$ выполнено $2p+1<2^p-1$ (\code{two\_mul\_add\_one\_lt\_mersenne}), так что $1<q<M_p$ и
$M_p$ составное. Оговорка $p>3$ честна: при $p=3$ спутник $q=7=M_3$ сам есть число Мерсенна и прост,
так что делитель совпал бы с делимым.
\end{proof}

\begin{theorem}[\code{mersenne\_composite\_examples}]\label{thm:mersenne_composite_examples}
\green{}\; $23\mid M_{11}=2047$, и $M_{11}$ составное.
\end{theorem}

Здесь $p=11\equiv3\pmod4$, $q=23$ просто, и $2047=23\cdot89$ распадается ровно на предсказанный
делитель. Через ограниченный фронт манифестации $3\pmod4$ это даёт, условно на открытой гипотезе
Софи Жермен, неограниченно много составных чисел Мерсенна
(\code{mersenneComposites\_unbounded\_of\_noEngine\_boundary\_and\_sgManifestation}, \green{},
условно) --- односторонний мост на \emph{составную} сторону того же $M_p$, на \emph{простую} сторону
которого граница Мерсенна ставить отказалась.

\subsection{Гольдбах и Лежандр: разрешимый свидетель}

Для Гольдбаха и Лежандра отклонение --- это объект-данное, вдобавок \emph{поточечно разрешимое} ---
сильнейшая форма непредъявимости в ряду. Каждый кандидат проверяется легче всего, а предъявить
подлинного свидетеля невозможно.

\begin{theorem}[\code{goldbachViolation\_ge\_52}]\label{thm:goldbachViolation_ge_52}
\green{}\; Всякий свидетель нарушения Гольдбаха сидит при $E\ge 52$.
\end{theorem}

\emph{Идея доказательства.} Нарушение Гольдбаха «$E\ge4$ чётно, не сумма двух простых» разрешимо
ограниченным перебором; \code{decide} проверяет \code{goldbach\_upTo\_52} (\green{}), что всякое
чётное $E<52$ разлагается, так что любая подделка ниже $52$ гибнет на месте. Литературная граница
$4\cdot10^{18}$ не формализована; зелен лишь $\ge52$. Закон манифестации квантифицирован по объекту
(без замка: якорь масштаба $M_0:=E$ есть само число), и суть проходит.

\begin{theorem}[\code{noGoldbachViolation\_of\_manifestation\_and\_boundary}]\label{thm:noGoldbachViolation_of_manifestation_and_boundary}
\green{}\; Нет двигателей, принятая граница и закон манифестации влекут отсутствие нарушений
Гольдбаха --- гипотеза Гольдбаха верна.
\end{theorem}

\emph{Идея доказательства.} Гипотетический свидетель $V$ даёт масштаб $M_0:=V.1$; граница поставляет
там разрешающую проекцию; закон превращает нарушение в бесконечное семейство потоков, из столкновения
которых собирается свидетель-двигатель --- убиваемый «нет двигателей». Через мост
\code{goldbachConjecture\_iff\_no\_violation} это доходит до самой гипотезы. Фронт Лежандра идентичен
(\code{noLegendreViolation\_of\_manifestation\_and\_boundary}, \green{}), но несёт единственного
безусловного спутника --- постулат Бертрана.

\begin{theorem}[\code{no\_desert\_doubles}]\label{thm:no_desert_doubles}
\green{}\; Простая пустыня не переживает удвоения: при $n\neq0$ между $n$ и $2n$ всегда есть простое.
\end{theorem}

\emph{Идея доказательства.} Прямая переупаковка постулата Бертрана
(\code{Nat.exists\_prime\_lt\_and\_le\_two\_mul}, \code{mathlib}~\cite{mathlib}, безусловно). Но он
\emph{не} решает Лежандра, и мы помечаем это машинно-проверяемо.

\begin{theorem}[\code{legendre\_interval\_shorter\_than\_bertrand}]\label{thm:legendre_interval_shorter_than_bertrand}
\green{}\; При $n\ge3$ выполнено $(n+1)^2<2n^2$: интервал Лежандра $(n^2,(n+1)^2)$ \emph{короче}
удвоения Бертрана $[n,2n]$.
\end{theorem}

Значит, зелёный Бертран бессилен на квадратичной трещине, и модуль фиксирует
\code{NoBertrandToLegendreImplicationClaimed} $=$ True. Зелёная теорема рядом с открытой задачей не
делает её зелёной; она размечает ровно то, где кончается доказанное.

\subsection{Совершенные числа: зелёный Евклид--Эйлер в обе стороны}

Обе стороны древнейшей теоремы \emph{Начал}~\cite{euclid-elements} теперь зелены, реэкспортированы из
архива \code{mathlib}~\cite{mathlib,wiedijk100}.

\begin{theorem}[\code{perfect\_of\_mersennePrime'}]\label{thm:perfect_of_mersennePrime}
\green{}\; Для $2\le p$ с простым $2^p-1$ число $2^{p-1}(2^p-1)$ совершенно --- направление Евклида
(\emph{Начала} IX.36).
\end{theorem}

\begin{theorem}[\code{evenPerfect\_eq}]\label{thm:evenPerfect_eq}
\green{}\; Всякое чётное совершенное число имеет вид $2^k(2^{k+1}-1)$ с простым множителем Мерсенна
--- обращение Эйлера.
\end{theorem}

Вместе эти две стрелки суть вся классическая теорема: простые Мерсенна и чётные совершенные числа ---
один объект. Это делает зелёной древнюю эквивалентность.

\begin{theorem}[\code{mersennePrimesInfinite\_iff\_evenPerfectUnbounded}]\label{thm:mersennePrimesInfinite_iff_evenPerfectUnbounded}
\green{}\; \code{MersennePrimesInfinite} эквивалентно неограниченности чётных совершенных чисел.
\end{theorem}

\emph{Идея доказательства.} Вперёд: конструкция Евклида с $N<p<2^p\le 2^{p-1}(2^p-1)$ гонит
совершенные числа вверх без потолка; назад: классификация Эйлера рассекает неограниченное чётное
совершенное число, а \code{Nat.prime\_of\_pow\_sub\_one\_prime} вынуждает растущий простой показатель.
Обе стороны открыты; зелена \emph{эквивалентность} --- древнейший вопрос \emph{Начал} и вопрос
Мерсенна буквально одно и то же. Чётная сторона конструктивна и заперта, поэтому отклонение ---
нечётное совершенное число --- может жить лишь на нечётной стороне; его свидетель поточечно разрешим,
и $\ge101$ зелено (литературная граница $>10^{2200}$ не формализована). Классическая форма Эйлера
гипотетического свидетеля машинно-проверена.

\begin{theorem}[\code{odd\_perfect\_euler\_form}]\label{thm:odd_perfect_euler_form}
\green{}\; Всякое нечётное совершенное $n$ имеет вид $n=q^{\alpha}m^2$ с простым $q$,
$q\equiv1\pmod4$, $\alpha\equiv1\pmod4$ и $q\nmid m$.
\end{theorem}

\emph{Идея доказательства.} Для нечётного $n$ множитель $2$ входит в $\sigma(n)=2n$ ровно в первой
степени; ввиду мультипликативности $\sigma$ чётным может быть ровно одно $\sigma(q^\alpha)$, что
вынуждает единственный нечётный показатель, а счёт по модулю $4$ закрепляет $q\equiv\alpha\equiv1
\pmod4$. Это форма Эйлера восемнадцатого века; вместе с «не менее трёх различных простых делителей»
она стягивает область несуществующего свидетеля --- не утверждая его существования ни так, ни иначе.

\subsection{Числа Ферма: вторая ставка против ожиданий}

Числа Ферма $F_k=2^{2^k}+1$ зеркалят Мерсенн на другой стороне решётки.

\begin{theorem}[\code{six\_mul\_fermatCenter}]\label{thm:six_mul_fermatCenter}
\green{}\; При $k\ge1$, с $c_k=(F_k+1)/6$, имеем $6c_k=F_k+1$, то есть $F_k=6c_k-1$: число Ферма
есть \emph{минус}-сторона своего центра.
\end{theorem}

\emph{Идея доказательства.} При $k\ge1$ показатель $2^k$ чётен, поэтому $F_k\equiv2\pmod3$ и, по
нечётности, $F_k\equiv5\pmod6$ (\code{fermatNumber\_mod\_six}); $F_0=3$ вне решётки и исключён. Цепь
центров здесь \emph{квадратична}, $c_{k+1}=6c_k^2-4c_k+1$ (\code{fermatCenter\_chain}, \green{}),
которую снова нельзя пилить --- шаг $3\to1$ не несёт собственного пилинга
(\code{isEmpty\_properCenterPeel\_three\_one}, \green{}), --- так что поддельного свидетеля нет.
Локализация свидетеля отсутствия --- рекорд программы.

\begin{theorem}[\code{fermatAbsenceBound\_ge\_65537}]\label{thm:fermatAbsenceBound_ge_65537}
\green{}\; Всякая граница отсутствия близнецов Ферма не меньше $65537$.
\end{theorem}

\emph{Идея доказательства.} Пара при $k=4$, а именно $F_4=65537$ и $F_4+2=65539$, оба простые,
существует зелено (\code{fermat\_twin\_instances}), и никакой свидетель не срежет её снизу; заметим,
что $k=3$ пары не даёт, $F_3+2=259=7\cdot37$ составное. Остальной фронт зеркалит Мерсенн слово в
слово, и цена точна.

\begin{theorem}[\code{fermatManifestationLaw\_iff\_unbounded\_of\_boundary}]\label{thm:fermatManifestationLaw_iff_unbounded_of_boundary}
\green{}\; При принятой границе закон манифестации эквивалентен неограниченности близнецов Ферма.
\end{theorem}

Декрет нёс бы ровно силу близнецов Ферма, но знак эвристики перевёрнут сильнее мерсеннова: сумма
$\sum 2^{-2^k}$ сходится почти мгновенно, простыми известны лишь $F_0$--$F_4$, тогда как
$F_5$--$F_{32}$ доказуемо составные, и консенсус вовсе не ждёт новых простых Ферма. Сам Ферма ставил,
что все $F_k$ просты, и проиграл эйлеровому $641\mid F_5$. Поэтому поле \code{fermatBoundary}
машинно-допустимо и, вновь, намеренно не взято. То, что для Мерсенна было разовым отказом по знаку
эвристики, с Ферма становится \emph{правилом} зоопарка: честная граница требует зелено-непредъявимого
свидетеля \emph{и} ставки, которой не стыдно.

\begin{remark}
Ни один из разделов выше не объявляет ни одной открытой задачи решённой.
\code{twin\_prime\_conjecture} остаётся \code{sorry}; Гольдбах, Лежандр, Софи Жермен, Полиньяк
(зазоры $4,6$), бесконечность простых Мерсенна и Ферма, а также существование нечётного совершенного
числа --- всё это остаётся \red{} открытым; таинт карантина неизменен. Зелена --- сжатая классическая
арифметика и честная разметка того, где начинается каждая ставка.
\end{remark}
